\section*{Summary}

A dark-matter particle that couples to the electron spin creates one magnon through the transverse spin, at leading order in $1/|\vect{S}|$ (\cref{sec:14.1}). In a two-sublattice antiferromagnet the one-magnon response carries the weight $w^{\mathrm{M}}_{\vect{q}} = |\vect{S}|(A_{\vect{q}} - B_{\vect{q}})/\omega^{\mathrm{M}}_{\vect{q}}$, which suppresses magnons near the zone center and enhances them near the ordering wavevector: 0.024 against 42 for magnons of \SI{5}{meV} in NiO. For NiO at $\sigmae = \SI{e-40}{cm^2}$ it opens near \SI{7.3}{keV} for a \SI{5}{meV} threshold and peaks at $3.2\times10^3$ events per kilogram-year near \SI{3}{MeV}. Near its edge the rate depends on the magnetic anisotropy, which the Heisenberg model leaves out. The spin-length sum rule estimates that linear spin-wave theory overstates the one-magnon weight by $\Delta S/|\vect{S}|$ at the next order, \SI{7.8}{\percent} for NiO, where the zero-point reduction $\Delta S = 0.078$ is that of four interleaved simple-cubic antiferromagnets.

The longitudinal spin creates magnon pairs (\cref{sec:14.2}). Their weight vanishes in a ferromagnet and at $\vect{q} = 0$, grows as $q^2$, and averages to 0.073 of the one-magnon weight in NiO, the part of the $\Delta S/|\vect{S}| = 0.078$ moved out of the one-magnon line that arrives in the continuum (the two differ by $\Delta S^2$). Because a pair shares its momentum, it reaches the energy $\frac12\mchi v^2$ at any mass, and the two-magnon channel opens at masses lower by $v_{\mathrm{max}}^2/4c_{\mathrm{M}}(v_{\mathrm{max}} - c_{\mathrm{M}})$, a factor 5.1 for NiO. Final-state interactions move the pair weight within the continuum without changing its integral; whether NiO has a bound pair is open. A magnon mixes with a phonon at linear order only through the dependence of the anisotropy energy on strain, which spin--orbit coupling and the dipolar interaction provide; the strain dependence of the exchange couples a phonon only to magnon pairs. The hybrids carry an off-diagonal spin--nuclear response that saturates the interference bound of \cref{sec:8.3}.

A boson absorbed at $\vect{k} \to 0$ drives the uniform precession (\cref{sec:14.3}). When the exchange is invariant under spin rotations, the uniform transverse spin response is one line at the Larmor frequency, with weight $2\langle S^z_{\mathrm{tot}}\rangle$, at any spin and temperature. The spin couples through $\lambda_S$, the fraction of the ionic angular momentum carried by the spin, which equals $g_J - 1$ only within a $J$ multiplet and is one for quenched $3d$ ions. The shape of the sample moves the line to the Kittel frequency and reweights its polarizations. For YIG at \SI{10}{\micro eV} and $g_{aee} = 10^{-13}$, with an assumed line width of $10^{-5}$ of the frequency, the peak rate is \SI{3.6}{s^{-1}.kg^{-1}}. An antiferromagnet absorbs at its resonance with the exchange-suppressed weight $|\vect{S}|\omega_{\mathrm{AFMR}}/(2\omega_E + \omega_A)$ per pair of sublattices, about $0.1\,|\vect{S}|$ in MnF$_2$, whose resonance lies at the measured \SI{1.1}{meV}.

Measured neutron spectra can replace a computed spin response, within limits (\cref{sec:14.4}). The magnetic cross-section at a given momentum transfer contains only the part of the spin response transverse to it, so without a model of the spin density no neutron measurement, polarized or not, gives the longitudinal element $\hat{\vect{q}}\cdot S\cdot\hat{\vect{q}}$. Within the localized-moment model, transverse data at several wavevectors $\vect{q} + \vect{G}$ fix the full tensor, and an isotropy hypothesis recovers it from two scattering angles. For quenched ions the neutron sees the moment with the spectroscopic $g$, the dark matter the spin, and the factor $[\lambda_S/(g_{\mathrm{ion}}/2)]^2$ converts one into the other, 0.83 for Ni$^{2+}$ at $g_{\mathrm{ion}} = 2.2$ and 0.73 at 2.34.

Targets without inversion or time-reversal symmetry add responses that a collinear magnet lacks (\cref{sec:14.5}). The current response of a chiral crystal acquires a part odd in $\vect{q}$, the finite-momentum optical activity, and the point group decides whether it exists; in $\alpha$-quartz the orientation-averaged gyration is only $-0.045$ of its value along the optic axis. Phonons of the lowest band of GaAs carry on average 0.12 units of angular momentum. Single-molecule magnets offer metastable lines, such as the line of Mn$_{12}$-acetate at \SI{1.098}{meV} in \SI{1}{T}, which keeps \SI{90}{\percent} of its weight below \SI{4.2}{K}. A multiferroic below its ordering temperature has a magnetoelectric cross response; near that temperature the response is open.

An ultralight field is a classical Gaussian wave, coherent for $t_{\mathrm{coh}} = 1/m_\phi v_0^2$, \SI{1.04}{s} at \SI{1}{neV} (\cref{sec:14.6}). Its average reproduces the golden-rule rate for records longer than the shorter of $t_{\mathrm{coh}}$ and the inverse width of the target's line, while a single record matches the ensemble only beyond $t_{\mathrm{coh}}$. A polarized nuclear-spin ensemble driven by the gradient or the dipole coupling of an axion-like field obeys the Bloch equations. Its transverse magnetization grows as $\Omega_{\mathrm{Rabi}}M_0t$ and saturates at $\Omega_{\mathrm{Rabi}}M_0T_2$, and a Gaussian drive turns the linear growth into a random walk once the record exceeds $t_{\mathrm{coh}}$. Each periodogram bin of signal plus noise is exponentially distributed. The resulting likelihood gives sensitivities that scale as $t_{\mathrm{int}}^{-1/2}$ below the coherence time and as $(t_{\mathrm{int}}t_{\mathrm{coh}})^{-1/4}$ above it, and it already contains the random amplitude, which a fixed-amplitude analysis misses by a factor 1.99 for the median limit and 2.73 at the detection threshold.

\Cref{tab:14-nuisances} collects the truncations and inputs of the chapter with their sizes. Near the lightest masses the anisotropy row dominates every other by more than an order of magnitude, and it is a model, not a calculation. For absorption into the uniform mode, the shape of the sample and the width of the line set the rate, and neither depends on the spin-wave expansion. In the coherent regime the largest effects are statistical: the length of the record against $t_{\mathrm{coh}}$ and $T_2$, and whether the likelihood contains the random amplitude.

\begin{table}[tbp]
  \caption{Uncertainty budget of the magnetic rates and of the coherent regime. Type in the classes of \cref{sec:8.8}: normalization, tilt (in $\omega$ or in the dark-matter mass), threshold and new support; ``unsized'' marks an item with no size known. Sizes are for NiO, YIG, MnF$_2$, GaAs, Mn$_{12}$-acetate and $^{207}$Pb in PMN-PT as named.}
  \label{tab:14-nuisances}
  \centering\scriptsize\setlength{\tabcolsep}{3pt}
  \begin{tabular}{@{}>{\raggedright\arraybackslash}p{0.27\linewidth}>{\raggedright\arraybackslash}p{0.13\linewidth}>{\raggedright\arraybackslash}p{0.48\linewidth}>{\raggedright\arraybackslash}p{0.08\linewidth}@{}}
    \toprule
    Truncation or input & Type & Size & Section \\
    \midrule
    One-magnon weight of linear spin-wave theory & normali\-zation & high by about $\Delta S/|\vect{S}|$: \SI{7.8}{\percent} (NiO), \SI{3.1}{\percent} for $|\vect{S}| = \frac52$, \SI{16}{\percent} for $|\vect{S}| = \frac12$ & \cref{sec:14.1} \\
    Two-magnon continuum & new support & 0.073 of the one-magnon weight over the zone (NiO); $w_\parallel/w_\perp < 0.05$ below \SI{0.5}{keV} & \cref{sec:14.2} \\
    Final-state interaction of the pair & tilt; new support & about 0.04 of the spin response, moved in energy (NiO); bound pair unsized & \cref{sec:14.2} \\
    Magnetic anisotropy near the gap & threshold; tilt in $\mchi$ & NiO, \SI{5}{meV} threshold, one-constant gap \SI{4.54}{meV}: rate $\times29$ at \SI{10}{keV} (upper value), $\times1.35$ at \SI{30}{keV}, $< \SI{3}{\percent}$ above \SI{100}{keV} & \cref{sec:14.1} \\
    Exchange constants from DFT~$+\,U$ & threshold & $J_2$ of NiO 11 to \SI{27}{meV} against 19.01: magnon energies $\times0.58$ to 1.42 & \cref{sec:14.1} \\
    Magnon--phonon mixing & tilt & gap \SI{3.0}{\percent} of the crossing energy at \SI{2.83}{GHz} (YIG model); NiO unsized & \cref{sec:14.2} \\
    Thermal magnons & new support & $\ee^{-\omega_{\mathrm{gap}}/T}$, exponent $-5\times10^3$ at \SI{10}{mK} (NiO) & \cref{sec:14.1} \\
    Sample shape & threshold; normali\-zation & YIG film within \SI{10}{\percent} of the Larmor line for $B_0 \ge 4.8\,\mu_0M_{\mathrm{s}}$; weight $\times1.74$ at \SI{10}{\micro eV}; saturated sphere exact & \cref{sec:14.3} \\
    Zero-point correction to the resonance weight & normali\-zation & 5 to \SI{23}{\percent} for $|\vect{S}| = \frac32$, $\omega_A/\omega_E = 2$ to 0.3; MnF$_2$ unsized & \cref{sec:14.3} \\
    Resonance of MnF$_2$ & threshold & fitted fields \SI{7}{\percent} below the measured \SI{1.1}{meV} and \SI{9.4}{T}; peak rate unchanged & \cref{sec:14.3} \\
    Width $\varGamma_{\mathrm{M}}$ of the uniform line & normali\-zation & peak rate $\propto 1/\varGamma_{\mathrm{M}}$; $10^{-5}\omega$ for YIG assumed & \cref{sec:14.3} \\
    Spin fraction $\lambda_S$ of quenched $3d$ ions & normali\-zation & $1 - O(\xi_{\mathrm{ion}}^2/\Delta_{\mathrm{CF}}^2)$ & \cref{sec:14.3} \\
    In-medium mixing of a dark photon & normali\-zation & set by the dielectric constant and the enclosure; unsized & \cref{sec:14.3} \\
    Neutron data for quenched ions & normali\-zation & $[\lambda_S/(g_{\mathrm{ion}}/2)]^2 = 0.83$ (Ni$^{2+}$, $g_{\mathrm{ion}} = 2.2$) or 0.73 (2.34); 0.973 at \SI{5.3}{keV} & \cref{sec:14.4} \\
    Longitudinal response from neutron data & new support & needs the localized-moment form or the isotropy hypothesis & \cref{sec:14.4} \\
    Chiral phonons as a band average & normali\-zation & GaAs 0.124, $\pm\SI{40}{\percent}$ with the dipole part of the force constants; rate-weighted within \SI{10}{\percent} only for $\mchi \ge \SI{2.2}{MeV}$; decay branching unsized & \cref{sec:14.5} \\
    Upper-well population of a molecular magnet & normali\-zation & Mn$_{12}$-acetate at \SI{1}{T}: line keeps \SI{90}{\percent} of $2|\vect{S}|$ below \SI{4.2}{K} & \cref{sec:14.5} \\
    Multiferroic response near $T_{\mathrm{CM}}$ & normali\-zation; tilt & unsized & \cref{sec:14.5} \\
    Record length against $t_{\mathrm{coh}}$ and $1/\varGamma_{\mathrm{tg}}$ & normali\-zation & mean within \SI{10}{\percent} of the golden rule for $t_{\mathrm{int}} \ge 13\,t_{\mathrm{coh}}$ ($T_2 \to \infty$) to $t_{\mathrm{coh}}$ ($T_2 = 0.1\,t_{\mathrm{coh}}$) & \cref{sec:14.6} \\
    Field amplitude fixed instead of Rayleigh & normali\-zation & limit too strong by 1.99 below $t_{\mathrm{coh}}$, 1.015 at $100\,t_{\mathrm{coh}}$; 2.73 at the detection threshold & \cref{sec:14.6} \\
    Linearized Bloch equation & normali\-zation & $(\Omega_{\mathrm{Rabi}}t)^2/6$ in the transient; \SI{30}{\percent} saturation at the PMN-PT threshold & \cref{sec:14.6} \\
    Inhomogeneous broadening; effective field $E^*$ & normali\-zation & $u_{\mathrm{sp}} = 3.8\times10^{-4}$; $E^*$ about \SI{50}{\percent} ($^{207}$Pb in PMN-PT) & \cref{sec:14.6} \\
    Line shape of the halo & tilt & standard halo only; unsized & \cref{sec:14.6} \\
    \bottomrule
  \end{tabular}
\end{table}
